%=== CHAPTER TWO ===
%% Rewritten 2026-09-18: adds inertial-parameter identification, aerial
%% manipulation and change detection, and restates the gap.

\chapter{Background and Related Work}
\label{ch:literature}
\begin{spacing}{1.5}
\setlength{\parskip}{0.22in}

\section{Model Predictive Control}

MPC repeatedly solves a finite-horizon optimal control problem, applies the first
input and shifts the horizon forward. Its practical value comes from combining
prediction, multivariable coordination and explicit constraints within one
optimisation \cite{rawlings2017mpc,mayne2014mpc}. Linear MPC is mature and
computationally efficient, but a single linear model is poorly suited to the
large attitude changes and thrust-direction coupling of a quadrotor. NMPC retains
the nonlinear dynamics and nonlinear tracking errors directly
\cite{grune2017nmpc}.

The price is computation. Multiple shooting, sequential quadratic programming,
Gauss--Newton Hessian approximations and real-time iteration schemes are widely
used to make NMPC practical \cite{diehl2005rti}. CasADi provides symbolic
differentiation and a high-level optimisation interface
\cite{andersson2019casadi}; acados generates efficient optimal-control solvers
around tailored SQP and QP components \cite{acados,verschueren2022acados}, with
HPIPM \cite{hpipm} as the quadratic subproblem solver.

\section{Quadrotor NMPC}

A quadrotor is differentially flat under common assumptions \cite{mellinger2011},
and many successful controllers exploit position and yaw references with a
cascaded inner attitude loop \cite{lee2010}. NMPC can instead optimise collective
thrust and body torque against the complete nonlinear rigid-body model, and this
formulation becomes advantageous as trajectory curvature and operating envelope
increase \cite{kamel2017linear,foehn2021}.

The controller in this dissertation predicts position, velocity, unit quaternion
and body rate. The attitude residual is the vector part of the relative
quaternion, which avoids treating the four quaternion components as independent
Euclidean coordinates. Input limits represent collective thrust and body-torque
authority. This formulation is nominally standard; the distinguishing issue is
that the equilibrium input and the plant parameters change at payload attachment
and again at release.

\section{Uncertainty and Adaptation in MPC}

Robust MPC protects against bounded uncertainty, while stochastic MPC describes
uncertainty probabilistically. Both can preserve margins without identifying the
current plant, but conservatism grows when the uncertainty set must cover both an
empty and a loaded vehicle. Adaptive MPC instead updates a parameterised model
online: certainty-equivalent schemes use the current estimate directly,
set-membership schemes maintain a feasible parameter set, and dual-control
schemes excite the plant to learn \cite{aswani2013learning,hewing2020learning}.

Learning-based MPC extends this using Gaussian processes, neural networks or
learned residual models. Data-driven MPC for quadrotors can compensate complex
aerodynamic effects \cite{torrente2021datadriven,torrente2021}, and neural
approximations can reduce online computation \cite{salzmann2023neural}; broader
surveys of learning for control under safety constraints are available
\cite{brunke}. Those methods are valuable when the uncertainty is poorly
structured. Payload pickup has stronger prior structure: rigid-body mechanics
already specifies how mass and attachment geometry enter the equations. The
present work therefore estimates a small physical parameter set rather than
learning a general residual map.

A distinct and influential alternative is to decline to identify the payload at
all, and to absorb its effect as a lumped matched disturbance. The
$\mathcal{L}_1$-adaptive augmentation of a quadrotor NMPC \cite{hanover2021} is
the representative of this family, and incremental nonlinear dynamic inversion
\cite{smeur} is its inner-loop counterpart. Such schemes are robust and require
no payload model, but their internal quantity is a force or an acceleration, not
a mass and a moment. They therefore cannot answer the question this dissertation
poses --- whether a payload is still attached --- because that question is not
represented in their state. The comparison is implemented here (an additive
lumped translational and rotational disturbance in the prediction model) but is
not among the frozen results.

\section{State and Parameter Estimation}

Recursive filters such as the extended Kalman filter propagate a local Gaussian
approximation. They are efficient and central to flight-state estimation, but
inequality constraints are not naturally hard constraints and abrupt parameter
changes violate the local linearisation assumption. MHE estimates a window of
past states and disturbances by solving a constrained optimisation; it can impose
bounds such as a physically admissible mass interval and can attribute residuals
using the same nonlinear model as the controller
\cite{rawlings2017mpc,rao2003mhe,rao2003}. Recent work learns the MHE weights
themselves \cite{neuromhe}, which is the principled version of the stage-weight
question that Chapter~\ref{ch:framework} settles negatively for this problem.

A moving window introduces a trade-off that this work encounters directly. A
longer horizon improves noise rejection, but the constraint that makes an
augmented parameter identifiable --- that it is constant across the window ---
becomes less defensible as the window spans more of the trajectory. The
implemented estimator uses 20 stages at \SI{0.1}{\second}.

\section{Inertial-Parameter Identification}

That the rigid-body equations are \emph{linear} in a particular set of inertial
parameters --- mass, first mass moment and the six inertia elements --- is a
classical result in manipulator identification \cite{atkeson1986}. The present
work identifies only the first of them online. For a payload grasped on the
thrust axis the first moment vanishes and the inertia increment reduces to a
known function of the mass, so the classical set collapses to one scalar; the
argument of Chapter~\ref{ch:model} is about which measurement channel should
carry that scalar on a vehicle whose rotational model is less trustworthy than
its translational one.

Payload-aware estimation for multirotors has previously used force--torque
sensing, external observers or explicit mass estimators \cite{tomic2017,
svacha2020, wuest2019, bohm2021}, and slung-load and aerial-manipulation work
treats the payload as a dynamic subsystem of its own \cite{sreenath2013,
thomas2014, ollero2022}. In the works reviewed here, whether a payload is
present is either assumed known or not represented as a quantity of its own.
This dissertation studies a setting in which the estimator is not told when a
payload is attached or released.

\section{Change Detection and Event-Triggered Estimation}

Deciding that a plant has changed, from a residual, is the subject of classical
change-detection theory \cite{basseville1993}, and event-triggered estimation is
an established response to it \cite{heemels2012}. The predecessor of the present
system used exactly this structure: a persistent step in the reconstructed
collective thrust triggered a de-weighting of the pre-event stages of the window.

Chapter~\ref{ch:framework} reports why that structure was retired for this
problem. A collective-thrust residual cannot distinguish a payload change from a
trajectory-mode change, a sustained vertical gust, or an aggressive manoeuvre,
because all of them move the same scalar. Any decision resting on that signal
alone inherits the ambiguity. The deployed design does not detect the change at
all: it waits for the estimated mass itself to persist at the bare-airframe value,
and accepts both the delay and the dependence on mass accuracy that this implies.
A formal comparison against a CUSUM detector on the same signals, under a matched
false-alarm constraint, is not performed and is listed among the gaps.

\section{Positioning and Research Gap}\label{sec:lit:gap}

Table~\ref{tab:related} places representative works from the families above
along the dimensions that matter for this dissertation. The summaries are
necessarily simplified to the formulation each work emphasises, and a dash means
that the aspect is not the subject of that work rather than that the work fails
at it.

\begin{sidewaystable}
\centering
\caption[Representative related work compared]{Representative related work compared along the dimensions relevant to
payload grasp and release. Entries summarise the formulation emphasised in each
cited work; ``---'' means the aspect is not its subject.}
\label{tab:related}
\scriptsize
\setlength{\tabcolsep}{3pt}
\renewcommand{\arraystretch}{1.15}
\begin{tabularx}{\textheight}{@{}>{\raggedright\arraybackslash}p{0.10\textheight} >{\raggedright\arraybackslash}p{0.11\textheight} >{\raggedright\arraybackslash}p{0.10\textheight} >{\raggedright\arraybackslash}p{0.10\textheight} >{\raggedright\arraybackslash}p{0.08\textheight} >{\raggedright\arraybackslash}p{0.09\textheight} >{\raggedright\arraybackslash}p{0.09\textheight} >{\raggedright\arraybackslash}p{0.07\textheight} >{\raggedright\arraybackslash}X@{}}
\toprule
\textbf{Work} & \textbf{Payload representation} & \textbf{Signals used} &
\textbf{Estimated quantities} & \textbf{Attach / release known?} &
\textbf{Release confirmation} & \textbf{Controller} & \textbf{Evidence} &
\textbf{Principal limitation for this problem} \\
\midrule
$\mathcal{L}_1$-adaptive NMPC \cite{hanover2021} & lumped matched force and torque disturbance & state estimate & disturbance (adaptation law) & not required & --- & NMPC with $\mathcal{L}_1$ augmentation & flight experiments & no payload quantity; presence not represented \\
Adaptive INDI \cite{smeur} & lumped inner-loop disturbance & angular acceleration, actuator feedback & control effectiveness & not required & --- & incremental attitude control & flight experiments & attitude loop only; no payload quantity \\
NeuroMHE \cite{neuromhe} & external disturbance wrench & state measurements & disturbance wrench; learned MHE weights & not required & --- & disturbance-compensated tracking & simulation and flight & payload not separated from other disturbances \\
Online parameter estimation \cite{wuest2019,bohm2021} & vehicle mass, centre of mass and inertia & onboard inertial and actuator signals & mass, CoM, inertia (and actuator parameters) & parameters assumed slowly varying & --- & not the focus & flight data & no release event; not coupled to a predictive controller \\
IMU-based inertia estimation \cite{svacha2020} & rigid-body inertia & IMU and rotor inputs & inertia & parameters assumed constant & --- & not the focus & experiments & inertia only; no grasp or release \\
External wrench estimation \cite{tomic2017} & external wrench & inertial and rotor signals & external force and torque; collisions & not required & contact events from the wrench & reflex reaction & flight experiments & payload not represented as mass and inertia \\
Aerial grasping \cite{mellinger2011} & rigid grasped object & external motion capture & object parameters after grasp & grasp commanded and known & --- & cascaded tracking control & flight experiments & attachment assumed known \\
This dissertation & centred rigid point mass on a nominal vertical arm & rotor speeds, odometry & composite mass; inertia derived & not given to the estimator & command-armed, estimator-confirmed detachment; force-balance confirmation with loss watchdog & NMPC over PX4 rate loop & software-in-the-loop only & mass collapses in manoeuvre; thrust-map sensitive; no hardware \\
\bottomrule
\end{tabularx}
\end{sidewaystable}

Within the works reviewed here, three aspects receive comparatively little
attention, and they define the scope of this dissertation.

First, \textbf{a payload model that does not need an event to switch off.}
Payload-adaptive formulations commonly estimate parameters conditioned on
attachment, so that some event must tell the model when the payload has gone.
A model whose payload-dependent terms are continuous functions of an estimated
quantity, and vanish as that quantity returns to the bare airframe, does not
need such an event for its dynamics.

Second, \textbf{the difference between a release command and a physical
detachment.} Systems that grasp and release usually treat the gripper command,
or a contact sensor, as the moment the payload leaves, and the mismatch between
the two is rarely measured. This dissertation uses the estimate to separate the
two once a release has been commanded, and measures both the delay this costs
and what happens when the same evidence is consulted before any command.

Third, \textbf{an explicit account of the costs.} Confirming a release on the
estimate delays it; estimating an absolute mass from an absolute thrust makes the
interface sensitive to a calibration constant; and resting every lifecycle
decision on one estimate makes the interface only as reliable as that estimate in
manoeuvre. Chapter~\ref{ch:evaluation} reports these costs rather than asserting
them to be small.

The dissertation consequently treats estimation and control as one architecture
rather than two algorithms. Estimation is judged by whether its outputs are
identifiable and useful for control; control is judged by whether it remains
consistent with the quantities being estimated, and by whether it knows when it
does not know.

\end{spacing}
\newpage
